\section{问题背景与问题重述}
\subsection{问题背景}

随着我国小区社群增加和用电需求的提高，光伏发电等分布式能源不断介入，复杂的能源网络给供电调节带来了严峻挑战。光伏供电和小区负载在时间上的峰值错位，使传统供电调节方式难以承担其任。为了解决分布能源本地消化，减少并网问题的发生，非孤岛式微网，一种能够整合小区负载，本地供电和储能设备的小型用电系统，成为了关键破局技术。

非孤岛式微网可通过与外部电网进行低价购电，高价出电的方式获取经济效益，降低购电成本。但是，微网调控面临着一个关键的阻碍：外网电价分时波动，光伏发电效率难以精确把控，同时储电设备存在容量上和充放电量上的双重约束。这使得在保证小区负载供电的条件下将购电支出最小化成为一个动态优化决策问题，它受到多方面约束的同时，在计划购电与实际出现偏差时，微网还会产生高价的紧急购电费和违约费，使得这个动态优化决策问题具有较高的复杂度。

在此环境下，迫切需要建立一套科学的、利于调控的为微网与外部电网电力调控策略，实现在电价波动，受环境影响的预测误差等情况下的动态调整，制定优秀的每日购电策略，从而减少购电成本，解决分布式电源的消纳难题，为能源环保提供保障。

\subsection{问题重述}

围绕微网与外部电网的购电调控，题目由确定性调度逐步扩展到预测不确定性、日内信息更新和实时电价波动四个层次。

\textbf{问题一：} 在每日电价和小区负载相同、已给出一天光伏发电功率预测的条件下，于每天 0:00 制定全天计划购电策略。要求任意时段供电不少于负载，并满足 0:00 与 24:00 储电量相同，在此基础上给出计划购电量、充放电量及购电费用。

\textbf{问题二：} 在每天电价规律相同而负荷、光伏随日期变化的条件下，于每天 0:00 提前制定全天计划购电量。若实际供电不足，则以交易时刻电价的 5 倍紧急购电，因此需要在正常计划购电成本与紧急补救风险之间进行权衡。

\textbf{问题三：} 在问题二基础上，每天 0:00、6:00、12:00、18:00 均可获得未来 24 h 的新光伏预报，并允许对尚未执行时段调整计划购电量。调增部分按 1.5 倍电价结算，调减部分承担 0.5 倍电价违约成本，需要建立利用新预报进行滚动修正的购电策略，并分析增加预报时点是否有必要。

\textbf{问题四：} 进一步考虑外部电网实时电价本身也随时间波动。在未来真实电价不可提前获知的条件下，需要对问题二和问题三的模型增加电价预测和价格不确定性处理，重新得到波动电价环境下的计划购电与滚动调整策略。

四个问题的逻辑关系可概括为：

\[
\boxed{
\text{问题一：确定性调度}
\rightarrow
\text{问题二：预测误差下的日前随机决策}
\rightarrow
\text{问题三：多时点信息更新下的日内再决策}
\rightarrow
\text{问题四：加入实时价格不确定性}
}
\]
